\section{Literature and conceptual framework}
\subsection{Infrastructure, incentives and adoption}
The infrastructure--adoption relationship is best understood as a sequence rather than a single outcome. Deployment changes the set of feasible journeys and can generate indirect network effects; adoption changes the population of vehicles and users; service records capture activity delivered by the available network. Research on network effects and subsidy design treats the first two elements as interdependent \citep{li2017,springel2021}, while work on policy complementarity shows that the translation from one to the next depends on the policy mix \citep{burra2024}.

The translation is conditional on instrument and setting. Chinese provincial evidence finds that vehicle subsidies, but not infrastructure subsidies, strengthen infrastructure--adoption links \citep{hu2025_charging}; adoption is consequently not a direct measure of charging delivered by supported facilities. Income, charger type and regional context also shape adoption responses \citep{gifford2026,gabehart2026}, and regional development studies describe industry and adoption trajectories rather than direct service activity \citep{wang2026_regional}.

\subsection{Charging activity and service}
Charging service contains several measurable layers. Utilization studies relate activity to charger type, venue, adoption and network conditions, whereas other work distinguishes occupancy, delivered energy and profitability under explicit cost assumptions \citep{borlaug2023,hecht2022}. Connection duration also combines charging with parking routines \citep{wolbertus2018}. These are not interchangeable proxies: each has a distinct denominator, temporal support and interpretation. A city-level activity-to-facility ratio is useful for comparing the observed network, but it does not become energy per port, occupancy or profitability by relabelling it.

Local composition is one source of this non-equivalence. Fast and slow chargers can respond differently to built-environment attributes \citep{zhao2024_utilization}; accessibility, urban functions and land use shape where and when stations are used \citep{du2024_accessibility,zheng2025_landuse}. Cities with similar facility indicators can thus record different activity because their equipment and locations differ. The city-level measure used here aggregates that heterogeneity deliberately: it describes recorded activity and intensity while leaving equipment, location and behavioural components as possible explanations rather than identified mechanisms.

Demand, successful charging and access can also diverge. Censored-demand and overflow models show that observed sessions may understate desired service when capacity binds, while habitual-use studies show that repeated routines can concentrate activity \citep{huttel2023,weekx2024,weekx2025_habits}. Planning reviews therefore distinguish activity served, charging opportunity, convenience and capacity \citep{patil2023}. Reliability and equity provide further dimensions that are not recovered from station counts or cost alone \citep{rempel2024,yu2025_reliability,zhangfan2025}. The empirical framework consequently preserves activity zeros and missing coverage as different source states and defines every ratio on its own compatible observation period.

\subsection{Operational responsibilities and evaluation}
The policy-mix framework distinguishes instruments, their interactions and the characteristics of the overall mix \citep{rogge2016}. For charging, this shifts attention from the number of policy documents or subsidies to the responsibilities assigned across infrastructure providers, operators and public authorities. Chinese research on charging subsidies and collaborative governance considers operating performance, feedback and administrative coordination \citep{hu2026,zhang2024}; international policy analysis similarly emphasizes utility coordination, data sharing and enforceable management arrangements \citep{tilly2025}. A continuing service is therefore an institutional arrangement as well as a physical stock.

Four non-exclusive functions make that arrangement observable in local texts. Information and supervision create visibility through platform connection, reporting or monitoring. Operation and maintenance concern the condition and running of equipment. Continuing-operation or payment provisions extend a duty over time or connect disbursement to later provision. Assessable or subsidy-linked conditions state criteria against which eligibility or performance can be checked. One clause may carry several functions: a payment rule can require continued operation, while a reporting rule can support verification without guaranteeing compliance. Coding the functions identifies the administrative object that a later evaluation could inspect; it does not measure implementation.

The conceptual link to the empirical outcomes is therefore one of alignment. Activity and facility indicators describe observed service and stock on compatible dates; policy coding describes formal expectations; annual associations and the governance timing comparison examine co-movement under their respective supports. None of these layers alone identifies a policy mechanism. Together they permit a differentiated assessment in which geographic coverage, activity growth, equipment condition, reporting visibility and continuity are treated as connected objectives with distinct evidence requirements.
